\chapter{Conclusão}
\label{chap:conclusao}

Esta dissertação testou se a população que vive numa bacia hidrográfica está associada ao microplástico medido na foz dessa bacia, na América do Norte e Central, dos Estados Unidos e sul do Canadá até o Canal do Panamá. Para isso, construiu um \textit{pipeline} que cruza três bases públicas e independentes: o relevo do GEBCO 2024, a população do GPW v4 e as contagens de \textit{pellets} plásticos do \textit{Nurdle Patrol} no banco de microplásticos marinhos da NOAA.

\section{Síntese dos resultados}

\textbf{Bacias.} O IPBA, com todos os pixels de terra vizinhos ao oceano como saídas, particionou o continente em 11.851 bacias costeiras sem nenhuma lista de rios. As 138 maiores cobrem 90\% da área de terra e reproduzem as grandes bacias conhecidas, do Mississippi ao Balsas, e somam 74\% da população da região.

\textbf{Manchas urbanas.} A calibração do CCA pela lei de Zipf, feita numa grade de 15 mil combinações de parâmetros por um algoritmo de união e busca que reproduz o CCA exatamente, escolheu $D^{*} = 1.010$ hab/km$^2$ e $\ell = 1{,}7$ km. Com esses valores, as 12.185 manchas urbanas seguem a lei de Zipf no topo do \textit{ranking} ($\hat{\alpha}$ entre 0,93 e 1,08 para as 50, 100 e 200 maiores), mas a distribuição completa é côncava, com $\alpha \approx 1{,}22$ na cauda ajustada pelo método CSN. O piso físico para $\ell$ mostrou-se necessário para evitar soluções degeneradas.

\textbf{Repartição.} As 618 manchas que cruzam divisores de água somam 41\% da população urbana, e a sua repartição entre as bacias muda substancialmente a população atribuída a cada foz. O método atribuiu corretamente cidades interiores, como a Cidade do México, Guadalajara e Monterrey, a fozes distantes centenas de quilômetros, mas deixou grandes cidades litorâneas em estuários fragmentados, como Nova York, fora das bacias principais.

\textbf{Microplástico e população.} Com o raio de aceitação que contém a maioria dos sítios de coleta ($R = 38$ km), 52 bacias têm medição, 29 com contagem positiva. Nessas 29, a concentração de \textit{pellets} cresce de forma sublinear com a população, com $\beta = 0{,}28$ (IC de 95\%: 0,11 a 0,54) e $r = 0{,}37$. A relação é estável entre 10 e 50 km, não depende de nenhuma bacia isolada, mas é fraca: a população explica cerca de 13\% da variância, e a correlação de postos cai para $\rho_s = 0{,}25$ ($p = 0{,}07$) quando as contagens zero são incluídas. O indicador de quantidade (concentração $\times$ área) dá $\beta = 0{,}85$ e $r = 0{,}66$, mas dois terços desse expoente vêm da correlação entre população e área das bacias. Com medição garantida, a até 1 km da foz, só três bacias estão disponíveis.

\section{Resposta à pergunta de pesquisa}

Os dados públicos disponíveis sustentam, no máximo, uma associação \textbf{positiva, sublinear e fraca} entre a população de uma bacia e a concentração de \textit{pellets} plásticos na sua foz. Eles não permitem afirmar que a população determina o microplástico que chega ao mar, nem estimar com precisão o expoente dessa relação. A principal limitação não está nos métodos, mas nos dados de plástico: as coletas se concentram no Golfo do México, raramente coincidem com fozes, têm em geral uma única amostra por sítio e medem um tipo de microplástico, os \textit{pellets}, cuja fonte é mais industrial do que populacional.

Esse resultado é, em si, relevante. Os modelos globais de emissão de plástico por rios usam a população da bacia como variável explicativa central \cite{lebreton2017,meijer2021}. A fraqueza da associação encontrada aqui com medições de campo independentes sugere cautela na extrapolação desses modelos e reforça a necessidade de medições sistemáticas nas fozes.

\section{Contribuições}

As contribuições do trabalho são:

\begin{itemize}
	\item um \textit{pipeline} reprodutível e aberto que integra relevo, população e microplástico numa análise por bacia, executável a partir de dois \textit{notebooks};
	\item o uso do IPBA com todos os \textit{sinks} costeiros e de um critério de cobertura de área para selecionar as bacias principais sem listas manuais nem limiares de área;
	\item um algoritmo de varredura do CCA em tempo quase linear por valor de $\ell$, exato em relação ao CCA, que torna viável explorar milhares de combinações de parâmetros sobre grades de dezenas de milhões de células;
	\item um critério de calibração do CCA pela lei de Zipf robusto a variações do tamanho da amostra e com um piso físico para a distância de ligação;
	\item a repartição das manchas urbanas entre bacias e a demonstração de que ela afeta quase metade da população urbana;
	\item uma análise explícita da sensibilidade ao raio de aceitação, que evita escolher o raio pelo resultado.
\end{itemize}

\section{Trabalhos futuros}

Os resultados sugerem as seguintes extensões.

\begin{itemize}
	\item \textbf{Agrupar \textit{sinks} de estuários.} Unir os \textit{sinks} de uma mesma baía ou estuário antes do IPBA, por exemplo agrupando pixels de costa separados por menos de alguns quilômetros de água, permitiria associar cidades como Nova York a uma única foz.
	\item \textbf{Relevo em resolução maior.} Uma implementação compilada do IPBA, com rótulos de 32 bits, permitiria usar o GEBCO na resolução nativa de 15$''$, localizando as fozes com precisão de algumas centenas de metros e aumentando o número de bacias com medição garantida.
	\item \textbf{Outras medidas de plástico.} Aplicar o mesmo \textit{pipeline} a medições de microplástico secundário na água ou no sedimento de estuários, ou a medições de carga fluvial, que testariam a hipótese de forma mais direta.
	\item \textbf{Modelos para as contagens.} Tratar as contagens zero com modelos de contagem, como regressões binomiais negativas com inflação de zeros, em vez de excluí-las do ajuste em log-log.
	\item \textbf{Covariáveis industriais.} Incluir a localização de fábricas de resina e terminais portuários como covariável, para separar o efeito da população do efeito da indústria no caso dos \textit{pellets}.
	\item \textbf{Transporte costeiro.} Ponderar os sítios por um modelo simples de deriva litorânea, em vez de usar apenas o sítio mais próximo da foz.
	\item \textbf{Outras regiões.} Aplicar o \textit{pipeline} a outras regiões com boa cobertura de coletas, como a Europa ou o litoral brasileiro, onde há programas de monitoramento de \textit{pellets} em praias.
\end{itemize}
